\section{A sharper dyadic estimate for the degenerate orbital integral}
\label{ods-section}

The last summation in the degenerate amplitude estimate is a
measure-theoretic argument involving two size functions and two
Gram determinants. We prove it independently and improve both its
constant and its decay exponent. The geometric estimates supplying
its hypotheses remain separate; in particular the relation between
the determinant and the size functions cannot be omitted.
\cite{alpo-uniform-source}

Let $(X,\nu)$ be a measure space. All integrals of nonnegative
functions below are allowed initially to be infinite. Let
$A,B,S,S',Y,Y'$ be measurable nonnegative real functions, with
$S,S'\ge1$ and $Y,Y'>0$ almost everywhere. Fix integers
$r\ge1$, $j\ge0$, real numbers $m\ge1$, $e\ge0$, and
\[
 q\ge e+3r+1,\quad H,\Delta,\sigma\ge1,
 \qquad D,K,K'>0.
                                                               \tag{OD1}
\]
Suppose that almost everywhere
\begin{align}
 YY'D&\le2^j,& D&\le(SS')^r,\nonumber\\
 Y,Y'&\le\max(S,S')^r\Delta,&
 \sigma&\le\max\left(1,\frac{(SS')^r}{(2H)^{2r}}\right).
                                                               \tag{OD2}
\end{align}
For $i\ge0$ write
\[
 E_i(S)=\{S\le2^{i+1}H\}\cap
       \bigl(\{i=0\}\cup\{2^iH<S\}\bigr),
\]
where $\{i=0\}$ means all of $X$ if $i=0$ and the empty set
otherwise. Define $E_i(S')$ in the same way. Suppose that for all
$i,i'\ge0$ and $y>0$,
\begin{align}
 \int_{E_i(S)\cap\{S'\le2^{i'+1}H\}\cap\{y<Y\le2y\}}
                    A^2\,d\nu
 &\le K(2^{i+i'+1}H)^e y^{2m}2^{-2qi},\label{ods-first}\\
 \int_{E_{i'}(S')\cap\{S\le2^{i+1}H\}\cap\{y<Y'\le2y\}}
                    B^2\,d\nu
 &\le K'(2^{i+i'+1}H)^e y^{2m}2^{-2qi'}.
                                                               \label{ods-second}
\end{align}

\begin{theorem}\label{ods-dyadic}
Under these hypotheses, $AB$ is integrable and
\begin{equation}\label{ods-bound}
 \int_X AB\,d\nu
 \le 4\,2^{-j}\sqrt{KK'}(2H)^{e+4r}\Delta^2
        \left(\frac{2^j}{D}\right)^m
        \sigma^{-(q-e-3r-1)/r}.
\end{equation}
No local compactness or finiteness assumption on $\nu$ is required.
\end{theorem}

\begin{proof}
Discard a null set so that the pointwise hypotheses hold everywhere.
The sets $E_i(S)$ form a disjoint countable partition, and likewise
for $S'$. On their intersection put
\[
 u=2^{i+i'},\qquad R=2Hu,\qquad M=R^r\Delta,
 \qquad N=2^j/D.
\]
Then
\begin{equation}\label{ods-size-product}
 \max(S,S')\le R,\qquad SS'\le u(2H)^2,
 \qquad \sigma\le u^r.
\end{equation}
In particular $0<Y,Y'\le M$. Subdivide into the measurable sets
$P_{i,i',a,b}$, for $a,b\ge0$, on which
\[
       M2^{-a-1}<Y\le M2^{-a},\qquad
       M2^{-b-1}<Y'\le M2^{-b}.
\]
These form a disjoint countable partition of the intersection.
Set $y=M2^{-a-1}$, $y'=M2^{-b-1}$ and $w=yy'$.
On a nonempty piece the product condition in (OD2) gives
$w\le N$. Conditions \eqref{ods-first}--\eqref{ods-second},
followed by Cauchy--Schwarz on this piece, give
\begin{equation}\label{ods-piece-cs}
 \int_{P_{i,i',a,b}}AB\,d\nu
 \le \sqrt{KK'}R^e w^m u^{-q}.
\end{equation}
The square integrals in this use of Cauchy--Schwarz are finite by
the same two hypotheses. Empty pieces contribute zero.

We retain two factors that can be lost in a rougher summation.
First,
\[
 w=M^2\,2^{-a-b-2}
     =(2H)^{2r}u^{2r}\Delta^2\,2^{-a-b-2}.
\]
Second, on each nonempty piece the product bound in
\eqref{ods-size-product} gives
\[
                    N^{-1}=D2^{-j}
                          \le2^{-j}u^r(2H)^{2r}.
\]
As $m\ge1$ and $w\le N$, these imply
\begin{equation}\label{ods-power}
 \begin{split}
 w^m&\le N^{m-1}w\\
 &\le N^m\,2^{-j-2}(2H)^{4r}u^{3r}\Delta^2\,2^{-a-b}.
 \end{split}
\end{equation}
Let $n=q-e-3r-1\ge0$. Since $\sigma\le u^r$, one has
$u^{-n}\le\sigma^{-n/r}$. Substitution in
\eqref{ods-piece-cs} yields
\[
 \int_{P_{i,i',a,b}}AB\,d\nu
 \le C\,2^{-i-i'-a-b},
 \quad
 C=2^{-j-2}\sqrt{KK'}(2H)^{e+4r}\Delta^2N^m\sigma^{-n/r}.
\]
Finally, monotone convergence and the four geometric series give
\[
 \int_XAB\,d\nu
 \le C\left(\sum_{h=0}^{\infty}2^{-h}\right)^4=16C.
\]
This is \eqref{ods-bound}; its right side is finite. The proof
also covers $n=0$, so equality is allowed in the bound for $q$.
\end{proof}

\begin{corollary}\label{ods-source-bound}
Suppose $m,e,q$ are integers, $m\ge1$, and
$q\ge e+4r+2$, with all the other hypotheses above. Then
\[
 \int_X AB\,d\nu
 \le16\sqrt{KK'}(2H)^{e+4r}\Delta^2
          (2^j/D)^m\sigma^{-(q-e-4r-1)/r}.
                                                               \tag{OD3}
\]
Thus the exact dyadic summation used for the submitted degenerate
orbital estimate follows. In fact \eqref{ods-bound} improves its
constant by a factor $4\cdot2^j$ and supplies one additional power
of $\sigma^{-1}$ under these same hypotheses.
\end{corollary}
\begin{proof}
The hypothesis on $q$ implies (OD1). In the decay exponent,
$q-e-3r-1=(q-e-4r-1)+r$. Since $\sigma\ge1$ and $j\ge0$,
the right side of \eqref{ods-bound} is at most the right side of
(OD3), with precisely the stated factors.
\end{proof}

\subsection{Why the determinant--size condition is essential}

\begin{example}\label{ods-counterexample}
The assertion (OD3) is false if the condition $D\le(SS')^r$ is
deleted while its other hypotheses are retained. To see this,
let $L\ge512$ be an integer and give
$X=\{0,1,\ldots,2L\}$ counting measure. Put
\[
 \begin{gathered}
 r=m=H=\Delta=\sigma=K=K'=1,
 \quad j=e=0,
 \quad q=6,
 \quad N=2^{-2L},\quad D=N^{-1},\\
 S(a)=S'(a)=1,\quad
 Y(a)=2^{-a},\quad Y'(a)=2^{a-2L},
 \quad A(a)=Y(a)/2,\quad B(a)=Y'(a)/2.
 \end{gathered}
\]
Here $YY'D=1$, both $Y,Y'$ are at most one, and the condition
on $\sigma$ holds. Each interval $(y,2y]$ contains at most one
of the distinct powers of two $Y(a)$. At that point
$A(a)^2=Y(a)^2/4\le y^2$. Thus \eqref{ods-first} holds when
$i=0$, independently of $i'$, and its domain is empty when
$i\ge1$. The identical argument with $Y'$ proves
\eqref{ods-second}. All the retained assumptions are satisfied.
But
\[
                  \int_XAB\,d\nu=(2L+1)N/4>256N,
\]
whereas the right side of (OD3) is exactly $256N$.
The failed hypothesis is $D\le1$.
This example is not a counterexample to the submitted theorem,
which includes the determinant--size condition explicitly.
\end{example}

\subsection{Role in the orbital estimate}

In the degenerate orbital application, $m=k-1-j$, $e=16r^3$,
$D$ is the absolute value of the nonzero leading determinant,
and $Y,Y'$ are the corresponding normalized Gram determinants
on the two sides. The functions $S,S'$ are the maxima of the
flag lengths and their inverses. The two one-sided square-integral
estimates supply \eqref{ods-first}--\eqref{ods-second}; their
constants retain the unipotent Haar normalization and the stabilizer
measure. All of these are actual hypotheses, not consequences of
dyadic summation. Theorem~\ref{ods-dyadic} proves the final
measure-theoretic passage as soon as those geometric and
representation-theoretic estimates have been supplied. It does
not assert that a polynomial tail estimate by itself controls a
degenerate orbital integral.


